% =====================================================================
%  CHƯƠNG 4 — HỆ THỐNG Y TẾ: ĐÁI THÁO ĐƯỜNG TÍP 2 TẠI VIỆT NAM
% =====================================================================
\chapter{Hệ thống y tế: bệnh đái tháo đường típ 2 tại Việt Nam}\label{chap:yte}

Chương này dựng phân hệ y tế của tác tử gồm ba lớp: kho tri thức văn bản về chẩn đoán, điều trị, dinh dưỡng;
\texttt{VietnamDiabetesKG} lưu các ràng buộc lâm sàng có nguồn; và một mạng MLP ước lượng nguy cơ đái tháo đường
từ hồ sơ xét nghiệm. Ba lớp được kiểm tra trên ba tình huống thường gặp ở người bệnh Việt Nam: uống rượu bia
khi đang dùng metformin, ăn thịt gà, và ăn cơm, bún, xôi.

\begin{phathien}{Lưu ý về phạm vi}
Hệ thống được xây cho mục đích học tập. Nội dung y khoa tổng hợp từ Quyết định 5481/QĐ-BYT, tiêu chuẩn của ADA và
các bảng chỉ số đường huyết đã công bố; mọi câu trả lời của tác tử đều kèm lời nhắc ``không thay thế chỉ định của
bác sĩ''.
\end{phathien}

\section{Bối cảnh dịch tễ và phác đồ của Bộ Y tế}
Đái tháo đường đang tăng nhanh tại Việt Nam cùng đô thị hóa, thay đổi chế độ ăn và giảm vận động. Theo Liên đoàn
Đái tháo đường Quốc tế, Việt Nam có khoảng 4 đến 5 triệu người trưởng thành mắc bệnh, tức khoảng 6 đến 7\% dân số
trưởng thành~\cite{idf2021atlas}, và các điều tra quốc gia ghi nhận tới khoảng một nửa số người mắc chưa biết mình
bị bệnh. Đái tháo đường típ 2 chiếm 90 đến 95\% số ca; cơ chế chính là đề kháng insulin ở cơ, gan và mô mỡ kèm suy
giảm dần chức năng tế bào beta của tụy.

Hướng dẫn chẩn đoán và điều trị đái tháo đường típ 2 hiện hành là tài liệu ban hành kèm Quyết định 5481/QĐ-BYT ngày
30/12/2020~\cite{byt5481}, thống nhất với tiêu chuẩn của Hiệp hội Đái tháo đường Hoa Kỳ~\cite{ada2024standards}.
Bảng~\ref{tab:chandoan} ghi bốn tiêu chuẩn chẩn đoán; đây chính là các thực thể loại \textsf{Biomarker} trong KG.

\begin{table}[H]\centering
\caption{Tiêu chuẩn chẩn đoán đái tháo đường và tiền đái tháo đường theo Quyết định 5481/QĐ-BYT}\label{tab:chandoan}
{\small\begin{tabular}{L{6.2cm}L{3.8cm}L{4.2cm}}
\toprule
Xét nghiệm & Đái tháo đường & Tiền đái tháo đường \\
\midrule
Glucose huyết tương lúc đói (nhịn ăn $\ge$ 8 giờ) & $\ge$ 126 mg/dL (7,0 mmol/L) & 100 -- 125 mg/dL \\
Glucose 2 giờ sau nghiệm pháp 75 g glucose & $\ge$ 200 mg/dL (11,1 mmol/L) & 140 -- 199 mg/dL \\
HbA1c (phòng xét nghiệm chuẩn hóa) & $\ge$ 6,5\% (48 mmol/mol) & 5,7 -- 6,4\% \\
Glucose bất kỳ kèm triệu chứng kinh điển & $\ge$ 200 mg/dL & -- \\
\bottomrule
\multicolumn{3}{l}{\footnotesize Không có triệu chứng kinh điển thì cần hai kết quả bất thường, cùng mẫu hoặc hai thời điểm.}
\end{tabular}}
\end{table}

Về điều trị, metformin là thuốc đầu tay; nhóm ức chế SGLT2 (empagliflozin) được ưu tiên khi người bệnh có suy tim
hoặc bệnh thận mạn; insulin bắt buộc với típ 1 và dùng khi típ 2 không kiểm soát được bằng thuốc uống. Mục tiêu
HbA1c cho đa số người trưởng thành là dưới 7\%. Thay đổi lối sống đi cùng mọi bậc điều trị: vận động ít nhất 150
phút mỗi tuần và giảm 5 đến 7\% cân nặng; ở người tiền đái tháo đường, chương trình DPP cho thấy can thiệp lối
sống giảm 58\% nguy cơ tiến triển thành bệnh~\cite{dpp2002}.

Ba đặc điểm của người bệnh Việt Nam khiến tri thức chung chưa đủ. Bữa ăn lấy \emph{cơm trắng} làm nền, chỉ số
đường huyết GI khoảng 73; các món xôi, bánh, bún, phở phổ biến; và \emph{rượu bia} gắn với tiệc tùng, đám cưới,
giỗ chạp, nên câu hỏi ``đang dùng thuốc thì uống được không'' là câu hỏi thật, thường xuyên. KG của bài vì vậy có
một phân hệ dinh dưỡng với các món ăn Việt và một phân hệ tương tác thuốc -- rượu.

\section{Thiết kế bản thể học và đồ thị tri thức y tế}
\subsection{Hệ thống thực thể y tế}
Bản thể học y tế có \SL{kg/health/classes} lớp (Hình~\ref{fig:ontology} ở Chương~\ref{chap:kg}), xếp quanh lớp gốc
\textsf{MedicalEntity}: bệnh, tình trạng và biến chứng, chỉ số sinh học, yếu tố nguy cơ, sáu nhóm thuốc, chất và
thực phẩm phân loại theo mức an toàn (\textsf{SafeFood}, \textsf{ModerateFood}, \textsf{RiskFood}), lối sống và
hướng dẫn chuyên môn. Hình~\ref{fig:schema} là lược đồ: các lớp chính và vị từ được phép nối chúng.

\hinhtikz{t4_schema}{Lược đồ \texttt{VietnamDiabetesKG}: các lớp chính và vị từ nối chúng theo ràng buộc miền
$\to$ miền giá trị của \texttt{health\_ontology()}}{schema}

\maNguon{kg_health__health_ontology}{Bản thể học y tế: \SL{kg/health/classes} lớp và \SL{kg/health/predicates}
vị từ kèm ràng buộc}

Mỗi thực thể có nhãn tiếng Việt, bí danh và thuộc tính lâm sàng: \textsf{HbA1c\_Cao} có \texttt{threshold=6.5},
\textsf{Suy\_Than\_Nang} có \texttt{egfr\_max=30}, \textsf{Com\_Trang} có \texttt{gi=73, gl=29}. Bảng~\ref{tab:kgtypes}
đếm thực thể theo lớp của cả hai KG, Hình~\ref{fig:kg_health_graph} vẽ toàn bộ đồ thị y tế.

\bang{tab_kg_node_types}{Số thực thể theo lớp của hai đồ thị tri thức}{kgtypes}

\hinhpdf{kg_health_graph}{1.0}{Toàn cảnh \texttt{VietnamDiabetesKG}: \SL{kg/health/nodes} thực thể xếp vòng theo
lớp, \SL{kg/health/edges} quan hệ}{kg_health_graph}

\subsection{Mối quan hệ tri thức và ràng buộc lâm sàng}
\texttt{VietnamDiabetesKG} có \SL{kg/health/predicates} vị từ, chia thành bốn phân hệ: \emph{chẩn đoán}
(\textsf{TIEU\_CHUAN\_CHAN\_DOAN}, \textsf{TIEN\_TRIEN\_THANH}, \textsf{YEU\_TO\_NGUY\_CO}), \emph{điều trị}
(\textsf{CHI\_DINH\_DAU\_TAY} và nghịch đảo \textsf{DIEU\_TRI}, \textsf{UU\_TIEN\_KHI},
\textsf{CHONG\_CHI\_DINH\_KHI}), \emph{an toàn thuốc} (\textsf{TUONG\_TAC\_NGUY\_HIEM}, \textsf{TANG\_NGUY\_CO},
\textsf{BIEN\_CHUNG\_CUA}, \textsf{XU\_TRI\_BANG}) và \emph{dinh dưỡng, lối sống} (\textsf{KHUYEN\_NGHI\_AN},
\textsf{HAN\_CHE\_AN}, \textsf{AN\_DINH\_LUONG}, \textsf{LA\_MOT\_DANG\_CUA} bắc cầu, \textsf{KHUYEN\_NGHI\_LOI\_SONG}).
Bảng~\ref{tab:kgrel} và Hình~\ref{fig:kg_relations} đếm số cạnh mỗi vị từ.

\bang{tab_kg_relations}{Số cạnh theo vị từ của hai đồ thị tri thức}{kgrel}
\hinhpdf{kg_relations}{1.0}{Phân bố quan hệ theo vị từ của hai KG}{kg_relations}

Đoạn mã~\ref{code:kg_health__build_health_kg} là toàn bộ dữ liệu của KG y tế. Mỗi lời gọi \texttt{F(\dots)} thêm
một bộ ba kèm \texttt{source} là mã tài liệu trong kho và các thuộc tính lâm sàng (mức độ, hậu quả, quy tắc, lý
do). Mọi bộ ba đều qua bộ kiểm tra ràng buộc của Chương~\ref{chap:kg} trước khi vào đồ thị.

\maDai{kg_health__build_health_kg}{Dữ liệu của \texttt{VietnamDiabetesKG}: thực thể có thuộc tính và bộ ba có nguồn}

\section{Mô hình học sâu chẩn đoán nguy cơ}
\subsection{Dữ liệu}
Công cụ \texttt{chan\_doan\_nguy\_co} của tác tử là một mạng MLP huấn luyện trên bộ dữ liệu dự báo đái tháo
đường~\cite{mustafa2023diabetes}: hồ sơ gồm giới tính, tuổi, tăng huyết áp, bệnh tim, tiền sử hút thuốc, BMI, HbA1c,
glucose máu và nhãn mắc bệnh. Sau khi bỏ trùng còn \SL{diabetes/data/n} hồ sơ, tỉ lệ mắc bệnh
\SL{diabetes/data/pos_ratePct}\%, chia phân tầng thành ba tập (Bảng~\ref{tab:diadata}). Hình~\ref{fig:dia_dist}
cho thấy HbA1c và glucose tách hai lớp rõ nhất, còn tuổi và BMI chồng lấn nhiều.

\bang{tab_dia_data}{Bộ dữ liệu chẩn đoán đái tháo đường sau tiền xử lý}{diadata}
\hinhpdf{dia_dist}{1.0}{Phân bố các đặc trưng chính theo nhãn mắc bệnh}{dia_dist}

Hồ sơ được mã hóa thành \SL{diabetes/data/n_features} đặc trưng: bốn cột số (tuổi, BMI, HbA1c, glucose), giới
tính, hai cờ bệnh nền và sáu cột one-hot của tiền sử hút thuốc; mọi đặc trưng được chuẩn hóa z-score theo trung
bình và độ lệch chuẩn của tập huấn luyện.

\maNguon{data__diabetes_features}{Mã hóa hồ sơ thành \SL{diabetes/data/n_features} đặc trưng}

\subsection{Kiến trúc và huấn luyện}
Mạng \texttt{DiabetesMLP} có hai tầng ẩn 128 và 64 nơ-ron (Hình~\ref{fig:mlp3d}). Với $\mathbf{x} \in
\mathbb{R}^{13}$:
\begin{align}
  \mathbf{h}_1 &= \operatorname{Dropout}_{0{,}2}\bigl(\operatorname{ReLU}(\operatorname{BN}(W_1\mathbf{x} + \mathbf{b}_1))\bigr) \in \mathbb{R}^{128},\\
  \mathbf{h}_2 &= \operatorname{Dropout}_{0{,}2}\bigl(\operatorname{ReLU}(\operatorname{BN}(W_2\mathbf{h}_1 + \mathbf{b}_2))\bigr) \in \mathbb{R}^{64},\\
  z &= \mathbf{w}_3^\top \mathbf{h}_2 + b_3, \qquad \hat{p} = \sigma(z).
\end{align}
Tổng cộng \SL{diabetes/mlp/params} tham số. Vì lớp dương chỉ chiếm chưa tới một phần mười, hàm mất mát là
entropy chéo nhị phân có trọng số lớp dương $w_+ = N_-/N_+$:
\begin{equation}
  \mathcal{L} = -\frac{1}{N}\sum_{i=1}^{N}\Bigl[w_+\, y_i \log \sigma(z_i) + (1-y_i)\log\bigl(1-\sigma(z_i)\bigr)\Bigr].
\end{equation}
Mạng được tối ưu bằng AdamW (tốc độ học $10^{-3}$, suy giảm trọng số $10^{-4}$), lô 1024 mẫu, giảm tốc độ học một
nửa khi PR-AUC trên tập kiểm định không tăng sau 3 epoch, dừng sớm nếu không cải thiện sau 8 epoch, tối đa 60
epoch. Huấn luyện chạy trên GPU trong \SL{diabetes/mlp/train_s} s.

\hinhtikz{t4_mlp3d}{Kiến trúc \texttt{DiabetesMLP}: 13 đặc trưng $\to$ 128 $\to$ 64 $\to$ 1 logit, mỗi tầng ẩn
gồm BatchNorm, ReLU và Dropout}{mlp3d}

\maNguon{dl_models__DiabetesMLP}{Mạng \texttt{DiabetesMLP} trong PyTorch}
\maDai{train_diabetes__fit_mlp}{Vòng huấn luyện MLP: BCE có \texttt{pos\_weight}, AdamW, giảm tốc độ học và dừng
sớm theo PR-AUC}

Đường học (Hình~\ref{fig:dia_history}) hội tụ êm: mất mát huấn luyện và kiểm định giảm cùng nhịp, không có dấu
hiệu quá khớp; mạng chạy đủ \SL{diabetes/mlp/epochs} epoch mà không kích hoạt dừng sớm, tốc độ học được giảm
nhiều lần ở nửa sau khi PR-AUC kiểm định chững lại.

\hinhpdf{dia_history}{1.0}{Đường học của MLP: mất mát BCE và chỉ số trên tập kiểm định qua các epoch}{dia_history}

\subsection{Kết quả và phân tích}
Ngưỡng quyết định được chọn để F1 lớn nhất trên tập kiểm định, ra \SL{diabetes/mlp/threshold}. Ngưỡng cao như vậy
là hệ quả trực tiếp của $w_+$: mất mát có trọng số đẩy mọi xác suất lên, nên 0,5 không còn là điểm cân bằng.
Trên tập kiểm thử, MLP đạt F1 \SL{diabetes/mlp/f1} và ROC-AUC \SL{diabetes/mlp/roc_auc}, vượt Logistic Regression
ở mọi chỉ số (Bảng~\ref{tab:diamodels}, Hình~\ref{fig:dia_compare}).

\bang{tab_dia_models}{MLP và Logistic Regression trên tập kiểm thử}{diamodels}

\hinhpdf{dia_roc_pr}{1.0}{Đường ROC và Precision -- Recall của hai mô hình trên tập kiểm thử}{dia_roc_pr}
\hinhpdf{dia_confusion}{1.0}{Ma trận nhầm lẫn tại ngưỡng đã chọn}{dia_confusion}
\hinhpdf{dia_compare}{1.0}{So sánh MLP và Logistic Regression theo từng chỉ số}{dia_compare}

Ma trận nhầm lẫn (Hình~\ref{fig:dia_confusion}) cho thấy tính cách của mô hình: chỉ
\SL{diabetes/mlp/confusion/0/1} ca dương tính giả trên \SL{diabetes/data/n_test} hồ sơ, precision
\SL{diabetes/mlp/precision}, nhưng bỏ sót \SL{diabetes/mlp/confusion/1/0} ca bệnh, recall
\SL{diabetes/mlp/recall}. Với một công cụ \emph{sàng lọc}, ngưỡng nên hạ thấp để đổi precision lấy recall; với một
công cụ trong tác tử trả lời người bệnh, ngưỡng cao giảm báo động giả. Bài giữ ngưỡng F1 và để tác tử trình bày
xác suất cùng ngưỡng thay vì một kết luận có/không.

Hai phân tích sau giải thích mô hình đã học gì. Độ quan trọng hoán vị (Hình~\ref{fig:dia_importance}) đo PR-AUC
giảm bao nhiêu khi xáo trộn một đặc trưng: HbA1c (\SL{diabetes/importance/HbA1c}) và glucose
(\SL{diabetes/importance/Glucose}) chiếm gần như toàn bộ, tuổi và BMI đóng góp nhỏ, các cờ còn lại gần bằng
không. Đường đáp ứng theo HbA1c (Hình~\ref{fig:dia_hba1c_curve}) cố định các đặc trưng khác và quét HbA1c: xác suất
vượt ngưỡng quyết định ở HbA1c \SL{dem/dia_cross_140}\% khi glucose 140 mg/dL và \SL{dem/dia_cross_100}\% khi glucose
100 mg/dL, sát ngưỡng chẩn đoán 6,5\% của Bộ Y tế. Mô hình tự học lại được tiêu chuẩn lâm sàng từ dữ liệu.

\hinhpdf{dia_importance}{1.0}{Độ quan trọng hoán vị của các đặc trưng}{dia_importance}
\hinhpdf{dia_hba1c_curve}{1.0}{Xác suất mắc bệnh theo HbA1c ở ba mức glucose}{dia_hba1c_curve}

Biểu đồ tin cậy (Hình~\ref{fig:dia_calibration}) cho thấy giới hạn quan trọng nhất: xác suất của MLP \emph{không
hiệu chuẩn}. Ở ngăn có xác suất dự đoán trung bình \SL{diabetes/calibration/5/pred}, tỉ lệ mắc thực tế chỉ
\SL{diabetes/calibration/5/obs}; chỉ ngăn trên cùng là khớp. Đây vẫn là hệ quả của $w_+$. Vì vậy con số mà tác tử
và giao diện hiển thị phải đọc như một \emph{điểm nguy cơ} so với ngưỡng, không phải xác suất mắc bệnh; muốn dùng
như xác suất cần hiệu chuẩn thêm bằng Platt scaling hoặc hồi quy isotonic trên tập kiểm định. Bộ dữ liệu cũng là
dữ liệu mô phỏng có quan hệ rất sắc giữa HbA1c và nhãn, nên hồ sơ có HbA1c trên 7\% gần như luôn được chấm sát
100\%; mô hình không nên dùng cho quyết định lâm sàng thật.

\hinhpdf{dia_calibration}{1.0}{Biểu đồ tin cậy: xác suất dự đoán bị đẩy cao so với tần suất thực tế}{dia_calibration}

Khi chạy trong tác tử, lớp \texttt{DiabetesRiskModel} nạp trọng số và tham số chuẩn hóa, điền giá trị mặc định
cho trường người dùng không cung cấp và trả mức nguy cơ ba bậc (Đoạn mã~\ref{code:dl_models__DiabetesRiskModel.predict}).

\maNguon{dl_models__DiabetesRiskModel.predict}{Suy luận nguy cơ từ hồ sơ trong tác tử}

\section{Phân tích các tình huống lâm sàng thực tế}
Ba bài toán dưới đây đều là câu hỏi người bệnh Việt Nam thường hỏi, và mỗi bài cần một loại suy luận khác nhau
trên KG. Bảng~\ref{tab:demohealth} là kết quả nguyên văn của các truy vấn KG tương ứng.

\bang{tab_demo_health}{Kết quả truy vấn \texttt{VietnamDiabetesKG} cho ba bài toán lâm sàng}{demohealth}

\subsection{Bài toán 1: tương tác giữa metformin và rượu bia}
Câu hỏi ``đang uống metformin, tối nay đi tiệc uống bia được không'' cần hai mẩu tri thức. Mẩu thứ nhất là một
cạnh trực tiếp $\textsf{Metformin} \xrightarrow{\textsf{TUONG\_TAC\_NGUY\_HIEM}} \textsf{Ruou\_Bia}$ với mức độ
\texttt{cao}. Mẩu thứ hai không được ghi tường minh mà là một \emph{đường hai bước}: cả metformin và rượu bia đều có
cạnh \textsf{TANG\_NGUY\_CO} tới \textsf{Nhiem\_Toan\_Lactic}. Hàm \texttt{drug\_warnings} lấy cạnh trực tiếp và
giao hai tập nguy cơ để tìm nguy cơ chung (Đoạn mã~\ref{code:kg_health__drug_warnings}). Cơ chế sinh học đằng sau
được vẽ ở Hình~\ref{fig:metformin}: ethanol tiêu thụ NAD$^+$ ở gan và ức chế tân tạo glucose từ lactat, metformin
cũng giảm tân tạo glucose, hai tác động cộng lại làm lactat ứ đọng.

\hinhtikz{t4_metformin}{Cơ chế tương tác metformin -- rượu bia: hai con đường cùng làm giảm sử dụng lactat ở gan,
dẫn tới nguy cơ nhiễm toan lactic}{metformin}

\maNguon{kg_health__drug_warnings}{Cảnh báo tương tác: cạnh trực tiếp và nguy cơ chung qua đường hai bước}

Câu trả lời đúng không phải ``tuyệt đối cấm'' mà là cảnh báo có mức độ: tránh uống quá mức và tránh say, không
uống khi đói, theo dõi dấu hiệu bất thường; với người dùng sulfonylurea hoặc insulin, rượu còn làm tăng mạnh nguy cơ
hạ glucose huyết tới 24 giờ sau khi uống. Khung dưới đây là vết suy luận thật của tác tử cho câu hỏi này, do
\texttt{build.py} sinh từ \texttt{agent.json}.

\begin{vetsuyluan}{Vết suy luận ReAct thật cho câu hỏi y tế (kịch bản 1)}
\input{generated/trace_health}
\end{vetsuyluan}

Bộ kiểm chứng đã từ chối bản nháp của LLM và trả bản tổng hợp theo khuôn. Bản nháp (dưới đây) đúng về nội dung,
nhưng không ghi số hiệu bằng chứng [n] nào, nên không truy vết được từng ý về nguồn; Mục~\ref{sec:agent-ketqua}
phân tích hiện tượng này trên toàn bộ kịch bản.

\begin{vetsuyluan}{Bản nháp của Qwen2.5-1.5B (bị bộ kiểm chứng từ chối vì thiếu trích dẫn)}
\input{generated/llmdraft_health}
\end{vetsuyluan}

\subsection{Bài toán 2: dinh dưỡng thịt gà}
``Người tiểu đường ăn thịt gà được không'' không có một câu trả lời: ức gà bỏ da luộc là đạm nạc 31 g/100 g, không
carbohydrate, được khuyến nghị; gà rán và da gà nhiều chất béo bão hòa, làm tăng đề kháng insulin, nên hạn chế.
Trong KG, \textsf{Thit\_Ga} là một \textsf{FoodGroup}, các dạng cụ thể nối tới nó bằng \textsf{LA\_MOT\_DANG\_CUA}.
Hàm \texttt{food\_advice} nhận ra câu hỏi về \emph{nhóm} thực phẩm và duyệt xuống các dạng con trước khi tra lời
khuyên (Đoạn mã~\ref{code:kg_health__food_advice}), nên trả cả hai chiều khuyến nghị thay vì một câu ``được'' hay
``không'' sai một nửa.

\maNguon{kg_health__food_advice}{Lời khuyên dinh dưỡng: duyệt nhóm thực phẩm xuống các dạng cụ thể}

\subsection{Bài toán 3: định lượng tinh bột và ẩm thực truyền thống}
Chỉ số đường huyết GI so sánh tốc độ một thực phẩm làm tăng glucose máu với glucose chuẩn (GI = 100); GI dưới 55
là thấp, trên 70 là cao~\cite{atkinson2008gi}. Tải lượng đường huyết GL $=$ GI $\times$ lượng carbohydrate (g) $/
100$ tính cả khẩu phần. Các món tinh bột Việt Nam trải khắp thang GI (Hình~\ref{fig:gi}): bún, phở GI khoảng 53, gạo
lứt 68, cơm trắng 73, khoai tây luộc 78, xôi nếp rất cao. KG không cấm cơm mà mã hóa ba mức: \textsf{AN\_DINH\_LUONG}
cho cơm trắng (tối đa khoảng một bát con mỗi bữa, ăn rau trước), \textsf{HAN\_CHE\_AN} cho khoai tây và xôi,
\textsf{KHUYEN\_NGHI\_AN} cho rau xanh (ít nhất nửa đĩa). Truy vấn SPARQL Q2 ở Chương~\ref{chap:kg} trả đúng danh
sách món có GI, sắp theo GI giảm dần, nhờ suy diễn kiểu từ các lớp \textsf{ModerateFood}, \textsf{RiskFood} lên
\textsf{Food}.

\hinhtikz{t4_gi}{Thang chỉ số đường huyết GI của các món trong KG và mô hình đĩa ăn: nửa đĩa rau, một phần tư
đạm, một phần tư tinh bột}{gi}

\subsection{Cơ chế phản ứng đa tầng trước tình huống phức tạp}
Câu hỏi thật hiếm khi chỉ một ý. Người bệnh nam 52 tuổi, HbA1c 7,3\%, glucose lúc đói 138 mg/dL, BMI 26,5 hỏi
``tôi đang uống metformin 500 mg, tối nay đi liên hoan có uống bia và ăn gà rán được không''. Tác tử phản ứng theo
ba tầng (Hình~\ref{fig:ui_health}):

\begin{enumerate}
  \item \textbf{Tầng học sâu}: hồ sơ xét nghiệm được gửi kèm, nên bộ lập kế hoạch gọi \texttt{chan\_doan\_nguy\_co}
        trước; MLP chấm điểm nguy cơ vượt ngưỡng, mức ``cao''.
  \item \textbf{Tầng đồ thị}: liên kết thực thể nhận ra \textsf{Metformin}, \textsf{Ruou\_Bia}, \textsf{Ga\_Ran};
        \texttt{kg\_y\_te} trả cảnh báo tương tác, các chống chỉ định của metformin và lời khuyên hạn chế gà rán.
  \item \textbf{Tầng văn bản}: \texttt{truy\_xuat\_rag} lấy các đoạn về rượu bia trong bữa tiệc và cơ chế nhiễm toan
        lactic từ tài liệu \texttt{h03-ruou-bia} để giải thích.
\end{enumerate}

Mỗi tầng trả lời một phần mà tầng khác không trả lời được: MLP không biết gì về bia, KG không biết người này có
nguy cơ cao, văn bản không chấm điểm được hồ sơ. Hình~\ref{fig:ui_health_food} là câu hỏi về cơm trắng trên cùng
giao diện, khi không gửi kèm hồ sơ: bộ lập kế hoạch bỏ qua MLP và chỉ dùng KG cùng RAG.

\anhUI{ui_health}{Trợ lý y tế trên Web App: hồ sơ xét nghiệm, vết suy luận ReAct, câu trả lời có trích dẫn và danh
sách bằng chứng}{ui_health}
\anhUI{ui_health_food}{Câu hỏi về cơm trắng không kèm hồ sơ: tác tử chỉ gọi KG và RAG}{ui_health_food}
